% =====================================================================
%  Lugworm: Complete Reference
%  Combines the former manual.md (quick reference) and pro-reference.md
%  (full technical reference) into a single document, organised so every
%  component has its compact parameter table immediately followed by
%  its full technical detail, one view per component instead of two
%  separate documents.
%
%  Compile with pdflatex (run twice for the table of contents / links):
%    pdflatex reference.tex
%    pdflatex reference.tex
% =====================================================================
\documentclass[11pt]{article}

\usepackage[a4paper,margin=2.2cm]{geometry}
\usepackage[utf8]{inputenc}
\usepackage[T1]{fontenc}
\usepackage{textcomp}
\usepackage{amsmath}
\usepackage{booktabs}
\usepackage{longtable}
\usepackage{array}
\usepackage{xcolor}
\usepackage{enumitem}
\usepackage{parskip}
\usepackage{titlesec}
\usepackage{fancyhdr}
\usepackage{hyperref}

\definecolor{lugwormblue}{RGB}{0,70,140}
\hypersetup{
  colorlinks=true,
  linkcolor=lugwormblue,
  citecolor=lugwormblue,
  urlcolor=lugwormblue,
  pdftitle={Lugworm Complete Reference},
  pdfauthor={Jasper Vollmert, IDAC, Hamburg University of Technology},
  bookmarksopen=true
}

\titleformat{\section}{\normalfont\Large\bfseries\color{lugwormblue}}{\thesection}{0.6em}{}
\titleformat{\subsection}{\normalfont\large\bfseries}{\thesubsection}{0.6em}{}
\titleformat{\subsubsection}{\normalfont\normalsize\bfseries}{\thesubsubsection}{0.6em}{}

\pagestyle{fancy}
\fancyhf{}
\fancyhead[L]{\small Lugworm: Complete Reference}
\fancyhead[R]{\small\leftmark}
\fancyfoot[C]{\thepage}
\renewcommand{\headrulewidth}{0.4pt}

\newcommand{\param}[1]{\texttt{#1}}
\newcommand{\iamdeg}{$^{\circ}$}
\newcommand{\heavy}{\textbf{[Heavy]}}

\setlist[itemize]{leftmargin=1.4em, itemsep=0.15em, topsep=0.2em}
\setlist[enumerate]{leftmargin=1.6em, itemsep=0.15em, topsep=0.2em}

\title{\color{lugwormblue}\textbf{Lugworm}\\[0.3em]\Large Complete Reference}
\author{Jasper Vollmert \\ \small Institute of Digital and Autonomous Construction (IDAC) \\ \small Hamburg University of Technology \\ \small jasper.vollmert@tuhh.de}
\date{\small Generated from \texttt{Lugworm} source (Components/*.cs, Core/*, docs/DECISION\_LOG.md): \today}

\begin{document}

\maketitle

\begin{abstract}
\noindent This document merges the quick-lookup manual and the full
technical reference for the \textbf{Lugworm} Grasshopper plugin
(assembly \param{Lugworm}, author \param{Lugworm}) into a single
view: every component below carries its compact parameter table
immediately followed by the technical detail behind it, so there is one
place to look rather than two documents to cross-reference. For a
plain-language introduction see \emph{beginner-guide.md}; for guided
practice see \emph{exercises.md}. This document assumes Grasshopper and
G-code familiarity.

The plugin documented here is the current, actively developed C\# tool.
An earlier Python GHPython prototype
(\param{ClaySlicer2}/\param{ClayMover2}/\param{ClayGcode2},
\param{clay\_core\_v2.py}) is superseded: per
\param{docs/DECISION\_LOG.md} entry D030, that Python toolkit is kept
only as a fixture-regeneration script for the test suite of this plugin, not
as a user-facing tool.
\end{abstract}

\tableofcontents
\newpage

% =====================================================================
\section{Overview and pipeline}
% =====================================================================

All components live under one Grasshopper tab, \textbf{Lugworm}, split
into seven numbered panels that match the order you would actually use
them in:

\begin{center}
\texttt{1 Slice\quad|\quad 2 Place\quad|\quad 3 G-code\quad|\quad 4 Robot\quad|\quad 5 Analyze\quad|\quad 6 Calibrate\quad|\quad 7 Print}
\end{center}

The core print pipeline:

\begin{verbatim}
your 3D geometry
      |
      v
+-------------+  contours (one    +-------------+  moved contours   +--------------+
|Slice |- branch/layer) ->|Place |------------------>|Lugworm G-code |--> .gcode
+-------------+                   +-------------+                   |   Generator  |
 "cut it into                      "put it on                       +--------------+
  layers"                           the bed,                         "write the machine
                                     fix the seam"                    instructions"
\end{verbatim}

From \param{G-code Generator}, the file either goes straight to a
Marlin/Duet printer, or through \param{G-code Restart} /
\param{Save G-code}. The machine itself is an explicit selection:
\param{Robot Definition} emits a profile string
(\texttt{iamlab-robot:v1:\ldots}) that every robot-dependent component
accepts on one \param{robot} input, and \param{Robot Simulation}
validates the program on the selected machine (see \S 6). Everything in
the core pipeline also runs with \textbf{no robot selected at all}.

Three further component groups sit alongside the pipeline:
\textbf{Analyze} (\param{Overhang}, \param{Print Stats})
for checking a slice or finished G-code before committing to print;
\textbf{Calibrate} (\param{DoE Design}, \param{Lugworm RSM
Evaluate}, \param{Desirability}, \param{Test Geometry})
for finding good settings for a new material via designed experiments
rather than guesswork; and \textbf{Print} (the \param{Lugworm Duet
\ldots} components) for direct Duet WiFi machine control, documented
in \emph{DUET\_CONTROL.md} rather than here.

\subsection{Component index}

\begin{longtable}{@{}p{3.9cm}p{1.8cm}p{2.0cm}p{6.9cm}@{}}
\toprule
\textbf{Display name} & \textbf{Nickname} & \textbf{Panel} & \textbf{Purpose} \\
\midrule
\endhead
\param{Slice} & Slice & 1 Slice & Mesh/Brep/Surface $\rightarrow$ contour tree \\
\param{Graph Mapper} & Graph Map & 1 Slice & Remap numbers through an editable curve \\
\param{Place} & Place & 2 Place & Bed placement and seam control \\
\param{Print Bed} & Bed & 2 Place & One bed definition for Place/Simulation/Live \\
\param{G-code Generator} & G-code & 3 G-code & Contour tree $\rightarrow$ Marlin/Duet G-code \\
\param{G-code Restart} & Restart & 3 G-code & Resume a failed print mid-way \\
\param{Save G-code} & Save & 3 G-code & Write G-code to a timestamped file \\
\param{Robot Definition} & Robot Def & 4 Robot & Select the machine profile (one wire) \\
\param{Robot Simulation} & Robot Sim & 4 Robot & Visualize/simulate/validate on the selected machine \\
\param{Overhang} & Overhang & 5 Analyze & Per-segment overhang/printability analysis \\
\param{Print Stats} & Stats & 5 Analyze & Length/volume/mass/time summary (robot-aware) \\
\param{DoE Design} & DoE & 6 Calibrate & Design-of-experiments run plan generator \\
\param{RSM Evaluate} & RSM & 6 Calibrate & Response-surface model fitting \\
\param{Desirability} & Desire & 6 Calibrate & Recommend settings from fitted models \\
\param{Test Geometry} & TestGeo & 6 Calibrate & Parametric calibration coupons \\
\param{Duet Connect/Status/Jog/Print/Live/Files/Terminal} & & 7 Print & Direct Duet WiFi machine control (see \emph{DUET\_CONTROL.md}) \\
\bottomrule
\end{longtable}

The pre-2.0 \param{Robot} component is hidden from the ribbon
but still deserializes and solves in saved documents; \param{Lugworm
Robot Simulation} replaces it (\S 6).

\heavy{} marks components with a manual compute button (cached per
component instance) rather than automatic recompute-on-change: \param{Lugworm
Slice}, \param{G-code Generator}, \param{G-code Restart},
\param{Overhang}, \param{RSM Evaluate}, and \param{Lugworm
Desirability}.

% =====================================================================
\section{Architecture and compute model}
% =====================================================================

\param{Lugworm/Components/*.cs} hold thin Grasshopper adapters
(parameter registration, caching, viewport preview); \param{Lugworm.Core/*}
holds the actual algorithms, organised by concern: \param{Analysis},
\param{Arcs}, \param{Clipping}, \param{Doe}, \param{Duet},
\param{Extrusion}, \param{Gcode}, \param{Geometry}, \param{Machine},
\param{Paths}, \param{Robot}, \param{Slicing}. This mirrors the Python
prototype's split (thin adapter + pure core) but as a compiled library
with a dedicated test project (\param{Lugworm.Core.Tests}, 980+
tests as of the 2.6.0 release).

\subsection{Heavy components and the compute gate}

\heavy{} components show a manual compute button instead of
recomputing on every canvas edit. Results are cached per component
instance (\param{SolveCache}, keyed by instance GUID) and evicted on
component deletion or document close. The robot simulator's cache key
is a content fingerprint (order-sensitive FNV-1a hash over every
sampled coordinate, extrude flags, and the full G-code string: a
reordering or interior-vertex edit is guaranteed to invalidate it,
fixed 2026-07-14 after an earlier bounding-box-sum fingerprint let stale
\param{danger=false} results survive a real edit). \param{Lugworm Save
G-code} and \param{G-code Restart} are not Heavy (cheap/live),
though \param{Save G-code} has its own manual \param{write} gate for a
different reason: avoiding re-writing a file on every recompute.

\subsection{Document units, language, and file I/O}

Version 2.6.0 introduced three behaviours that apply across the plugin:

\begin{itemize}
  \item \textbf{Document units.} Geometry is converted at the Rhino/Core
  boundary, so a document in millimetres, inches, or metres slices
  correctly and previews overlay the document geometry at document
  scale. Plain-number inputs documented in mm (\param{layerHeight},
  \param{width}, \param{spacing}, tolerances, \ldots) remain physical
  millimetres regardless of document units, and emitted G-code /
  Duet traffic is always plain mm. Changing document units
  mid-session invalidates Heavy caches (``no cached result: press
  Compute'') rather than replaying a stale-scale result.
  \item \textbf{Language/culture.} Machine-readable text (G-code,
  CSV, reports) is culture-invariant: a German-locale Windows
  emits decimal points, never commas. CSV \emph{inputs} (RSM tables,
  pump calibration) are detected per file as English (\texttt{1.5}) or
  German (\texttt{1,5}); a file mixing both conventions, using thousands
  separators, or with malformed quoting is refused with the offending
  line named, never silently misread.
  \item \textbf{Background file I/O.} File reads and writes (G-code
  save/restart, pump CSV, overhang CSV, RSM read/export) run in
  cancellable background workers so the canvas never
  blocks on a slow disk or network share. While a job is pending the
  component replays its last good outputs and reports
  ``writing/loading \ldots''; the result lands on the follow-up solve.
\end{itemize}

\subsection{Multi-geometry input}

\param{Slice}'s \param{geo} is a list parameter (single items
auto-promote). A single object emits the old flat \texttt{\{i\}} path
structure unchanged (backward compatible); a true list emits
\texttt{\{object;layer\}} paths. Downstream components (\param{Place},
\param{G-code Generator}, \param{Robot Simulation}) iterate branches generically and handle
both shapes. Multi-object printing is \textbf{sequential} (all of object
A's layers, then object B's, ...), order given by the input list or
overridden by the \param{printOrder} toggle-group input: deliberately
not layer-interleaved, since interleaving invites SCARA nozzle/part
collisions and travel time would dominate anyway; true interleaving
needs a collision-aware travel planner that does not exist yet.

\subsection{Geometry validation policy}

From \param{Core/Geometry} (\param{GeometryInput.Prepare} +
\param{MeshChecks.Analyze}):

\begin{itemize}
  \item \textbf{Rejected outright:} null/unsupported type, unmeshable Brep, zero faces, non-manifold edges.
  \item \textbf{Warned, still sliced:} open mesh (naked edges), degenerate faces, disconnected components, non-solid Brep.
  \item \textbf{Auto-fixed, reported:} inconsistent face normals via \param{Mesh.UnifyNormals}.
\end{itemize}

Breps still section via precise \param{BrepPlane} intersection
(unchanged from before), but are \emph{additionally} meshed for
validation/adaptive-slope/conical-shear use: a Brep pays a
meshing cost even in plain planar mode now, with no size-based opt-out
currently built.

% =====================================================================
\section{Slice: Panel 1 Slice \heavy{}}
% =====================================================================

\emph{``Slice a mesh/brep/surface (or a list of them) into a clay contour tree.''}

\begin{longtable}{@{}p{3.0cm}p{2.0cm}p{1.6cm}p{7.0cm}@{}}
\toprule
\textbf{Parameter} & \textbf{Unit} & \textbf{Default} & \textbf{Description} \\
\midrule
\endhead
\param{geo} & Geometry, list & required & Mesh/Brep/Surface(s) to slice; validated per \S 2.3. \\
\param{layerHeight} & mm & 3.0 & Layer spacing. \\
\param{firstLayer} & mm & 3.0 & \param{0} $\rightarrow$ half a normal layer. \\
\param{wallCount} & int & 1 & Number of perimeters. \\
\param{spacing} & mm & 10.0 & Wall/infill spacing: \textbf{hard error if $\leq 0$}. \\
\param{mode} & int & 0 & 0 Planar / 1 Conical. \\
\param{infill} & int & 0 & 0 None / 1 Zigzag / 2 Concentric / 3 Spiral / 4 Fermat. \\
\param{cornerRadius} & mm & 0.0 & Fillets sharp corners (over-fill correction). \\
\param{tol} & mm & 0.05 & Linearisation/simplification tolerance. \\
\bottomrule
\end{longtable}

\textbf{Toggle groups:}
\textbf{Adaptive} (\param{hMin} = layerHeight$\times$0.5, \param{hMax} =
layerHeight$\times$1.5, planar-only); \textbf{Conical} (\param{coneAngle} = 20\iamdeg, \param{nozzleClear} =
25\iamdeg, advisory collision warning only); \textbf{Zigzag} (\param{angle} = 0\iamdeg); \textbf{Spiralize} (menu toggle, no parameter: requires exactly one
closed contour per layer, all layers); \textbf{Meshing} (\param{meshDensity} = 0.5, \param{meshMaxEdge} = 0,
\param{meshMinEdge} = 0: controls Brep-to-mesh conversion quality for
validation/adaptive/conical); \textbf{Multi-object} (\param{printOrder}, an index list).

\textbf{Outputs:} \param{contours} (Curve tree), \param{layers} (int),
\param{info} (str: reports per-object issues; one bad object in a
list does not abort the whole solve).

\textbf{Non-planar specifics.} Because conical ring spacing follows the
sheared surface rather than a constant vertical Z step,
\textbf{extrusion volume in conical (and adaptive) mode is only
nominal} unless the \param{varh} group of the G-code component is
enabled (\S 5.1). The heat method (\param{mode = 2}) was
\textbf{retired in 4.0.0}: it emitted perimeter rings only and fell
back to planar on most real meshes. A saved definition carrying
\param{mode = 2} slices planar and reports a warning.

% =====================================================================
\section{Place: Panel 2 Place}
% =====================================================================

\emph{``Place contours on the bed and set the seam (clay / LDM).''} Not
Heavy (live).

\begin{longtable}{@{}p{3.0cm}p{2.0cm}p{2.2cm}p{6.4cm}@{}}
\toprule
\textbf{Parameter} & \textbf{Unit} & \textbf{Default} & \textbf{Description} \\
\midrule
\endhead
\param{contours} & Curve tree & required & From Slice. \\
\param{bed} & Bed definition & & Optional \param{Print Bed} wire; supplies the print centre when \param{printCenter} is unwired. \\
\param{printCenter} & Point, mm & (339, $-193$, 0) & The lab's SCARA bed ``sweet spot''; wins over the bed definition when wired. \\
\param{zOffset} & mm & 0.0 & Added after placement. \\
\param{move} & Vector & zero & Extra manual shift alongside centring. \\
\param{seamMode} & int & 1 & 0 Keep / 1 Align / 2 Stagger / 3 Nearest / 4 Guide / 5 Track-below. \\
\bottomrule
\end{longtable}

\textbf{Seam group:} \param{seamAngle} (0\iamdeg, modes 1 to 2),
\param{seamPoint} (mode 3), \param{seamCurve} (modes 4 to 5: ``Guide''
follows a supplied curve, ``Track-below'' keeps the seam aligned with
the layer underneath rather than a fixed target: new in the
2026-07-14 run).

Placement math: \(dx = printCenter.X - bboxCenter.X + move.X\) (and
analogously for Y, and Z with \param{zOffset}).

\textbf{Outputs:} \param{moved}, \param{seams} (point tree),
\param{info}. Placement is robot-independent: whether the placed part is
reachable is validated in one place only, \param{Robot
Simulation} (\S 6), so Place itself emits no reach warning.

\subsection{Print Bed}

\emph{``Define the print bed once: size, position and default print
center: and wire it to placement, simulation and live view.''} Not
Heavy. Emits a versioned definition string
(\texttt{iamlab-bed:v1:w=\ldots;d=\ldots;cx=\ldots}) consumed by
\param{Place}, \param{Robot Simulation} and \param{Lugworm
Live View}; an unrecognised string fails loudly rather than resolving to
wrong geometry.

Two separations are deliberate. \textbf{Plate centre is not print
centre}: \param{center} is the plate, \param{printCenter} is where a
part's bottom-centre lands: (468, $-441$) versus (339, $-193$) on the
lab machine. And \textbf{reach belongs to the arm, not the plate}: \param{reachMin} and \param{reachMax} of the bed default to ``ask the robot
profile'' and exist only so a physically restricted plate can narrow the
annulus. Every pre-existing bed input on the consuming components
survives and overrides the definition when wired, so saved documents are
untouched.

% =====================================================================
\section{G-code Generator: Panel 3 G-code \heavy{}}
% =====================================================================

\emph{``Convert any curve/polyline tree into material-extrusion G-code
(Marlin / Duet RRF).''}

\subsection{Core parameters}

\begin{longtable}{@{}p{3.4cm}p{1.8cm}p{1.6cm}p{6.5cm}@{}}
\toprule
\textbf{Parameter} & \textbf{Unit} & \textbf{Default} & \textbf{Description} \\
\midrule
\endhead
\param{crvs} & Curve tree & required & From Place. \\
\param{robot} & Robot definition & & Selects the coordinate frame (and machine defaults). \textbf{Unwired = legacy Potterbot frame}; see below. \\
\param{layerHeight} & mm & must match Slice & E cross-section height. \\
\param{width} & mm & 10.0 & Layer width. \\
\param{ramDiameter} & mm & 94.0 & Cartridge bore; warns (not a hard limit) outside 5 to 200\,mm. \\
\param{pressureCoeff} & & 1.2 & Overall flow multiplier, final calibration knob. \\
\param{printSpeed} & mm/s & 43 & Extruding feedrate. \\
\param{travelSpeed} & mm/s & 30 & Non-extruding feedrate. \\
\param{retractSpeed} & mm/s & 1 & Ram screw speed for retract \& prime (the screw moves in raw mm; retraction $\approx 0.1$\,mm, so keep this small). \\
\param{firstLayerSpeedFactor} & & 0.5 & Speed multiplier, first layer only. \\
\param{retraction} & 0/1 & & Master on/off for pull-back/prime behaviour. \\
\param{seamTol} & mm & & Seam-adjacent tolerance for coast/lead-in placement. \\
\param{firmware} & enum & Duet-RRF / Marlin & Changes exact G-code dialect emitted. \\
\param{relativeE} & bool & true & \texttt{M83}-style relative E. \\
\param{pressureOffset} & mm & & Plunger pull-back (``retraction'' for paste). \\
\param{primeExtra} & mm & & Extra prime on restart. \\
\param{coast} & mm & & Stop extruding this far before a path end. \\
\param{zHop} & mm & & Travel lift. \\
\param{hopSpeed} & mm/s & & Z-hop feedrate. \\
\param{footerZLift} & mm & & End-of-print lift. \\
\param{dwell} & ms & & Pause after prime. \\
\param{leadIn} & mm & & Flow ramp-in length. \\
\param{leadInBoost} & & & Flow multiplier during lead-in. \\
\param{minLayerTime} & s & & Slow a layer that would print faster than this. \\
\param{zOffset} & mm & 0.0 & Nozzle-to-bed calibration shim, added on top of geometry the slicer has already placed at its true height. The slicer owns layer height; leave at 0 unless the Z zero of the machine needs a correction. \\
\param{tol} & mm & & Curve simplification tolerance. \\
\param{savePath} & path, optional & & Write directly from this component (background write; the saved path is reported on the completion re-solve). \\
\bottomrule
\end{longtable}

\textbf{Coordinate frame.} \param{GcodeConfig.CoordinateFrame} comes from
the wired robot definition and is applied in the emitters (arc sense
flips with the reflection). An \textbf{unwired} \param{robot} input
resolves to \param{potterbot-scara-v4} and its calibrated legacy frame
(a mirror about $y = -x$): every G-code emitted before 2.0.0 carried
that reflection, so defaulting to identity would have silently mirrored
the lab's saved definitions. A Cartesian user must wire a definition
(\param{generic-cartesian}, \param{generic-duet-cartesian}, or
\param{none}) to get the identity frame. The active frame is reported in
\param{info} (\texttt{frame:\ldots}).

\subsection{Optional toggle groups}

\begin{itemize}
  \item \textbf{advance / retractaccel}: firmware pressure-advance and acceleration overrides (the \texttt{M900}/\texttt{M204}-equivalent hand-off to firmware).
  \item \textbf{arcs}: built-in arc fitting (\texttt{G2}/\texttt{G3}) instead of linear-only output.
  \item \textbf{arcwelder}: post-process via the ArcWelder CLI instead of the built-in fitter; \textbf{depends on \param{arcs}} (auto-disabled if arcs is off). See \S 8.3 for a critical historical bug here.
  \item \textbf{skirt}: purge ring before the real print.
  \item \textbf{floors}: solid base layer count.
  \item \textbf{headeroverride}: replace the auto-generated header entirely.
  \item \textbf{pump}: screw pump support (any speed-controlled pump): \param{nozzleDiameter}, \param{pumpCalibCsv}, \param{slipFactor}. See \S 8.1: \textbf{this group only calculates a recommended pump setting; it does not alter E-values in the emitted G-code.}
  \item \textbf{varh} (variable layer-height flow): \param{varhMin} (0.25), \param{varhMax} (2.5): clamps the per-point flow multiplier derived from \param{VariableFlow.LocalHeights}. \textbf{Turn this on whenever slicing used adaptive or conical mode.}
  \item \textbf{ecomp} (slope-compensated flow): \param{eCompRefMul} (1.0 = no-op), \param{eCompRefDeg} (45\iamdeg): experimental extra-flow correction on inclined walls via \param{SlopeFlow.Multipliers}. \textbf{NEEDS\_INPUT: uncalibrated; leave off until an inclined-wall calibration print has been run.}
\end{itemize}

\textbf{Outputs:} \param{gcode} (list), \param{gcode\_text} (str),
\param{preview}/\param{travels} (curve trees), \param{length} (mm),
\param{volume} (cm$^3$), \param{time} (min), \param{info}, \param{pump}
(text summary, populated only if the pump group is active),
\param{start} (first point), \param{startGcode} (a jog-to-start block,
useful for manually positioning before a resumed/robot print).

% =====================================================================
\subsection{G-code Restart}
% =====================================================================

\emph{``Resume a failed print mid-way, cutting the G-code at a Z height,
Duet \% or XYZ position and rebuild it from there.''} Heavy.

\begin{longtable}{@{}p{3.4cm}p{2.0cm}p{7.9cm}@{}}
\toprule
\textbf{Parameter} & \textbf{Unit} & \textbf{Description} \\
\midrule
\endhead
\param{gcode} & Text list & From the G-code component. \\
\param{z} / \param{percent} / \param{point} & mixed & Mutually exclusive cut selectors; priority \textbf{point $>$ percent $>$ z} if more than one is supplied. \param{point} is a Rhino document-coordinate XYZ point (for example, a Live or Console position output), not raw Duet machine XYZ. \\
\param{primeE} & mm & Re-prime distance at the resume point. \\
\param{travelSpeed}, \param{hopSpeed}, \param{retractSpeed} & mm/s & Resume-move feedrates. \\
\param{clearance} & mm & Safety lift before the resume travel. \\
\param{savePath} & path & Output location. \\
\bottomrule
\end{longtable}

\textbf{Outputs:} \param{gcode}, \param{gcode\_text}, \param{resume}
(the point execution resumes from), \param{printed}/\param{remaining}
(curve lists: useful to visually confirm the cut point before
committing), \param{info}.

% =====================================================================
\subsection{Save G-code}
% =====================================================================

\emph{``Write G-code to a folder with a date/time-stamped filename.''}

\begin{longtable}{@{}p{3.4cm}p{2.0cm}p{7.9cm}@{}}
\toprule
\textbf{Parameter} & \textbf{Unit} & \textbf{Description} \\
\midrule
\endhead
\param{gcode} & Text list & Input G-code. \\
\param{folder} & path & Destination folder. \\
\param{prefix} & str & Default \texttt{"print"}. \\
\param{write} & bool & Wire a \textbf{Button}, not a persistent Toggle, so it fires once per click rather than on every solve. \\
\param{dateTime} / \param{format} & & Override group; \param{format} default \texttt{"yyyy-MM-dd\_HH-mm"}. \\
\bottomrule
\end{longtable}

\textbf{Outputs:} \param{path}, \param{written} (bool: true only on
the solve where the write completed, not persistently), \param{info}.
The write itself runs in a background worker (2.6.0): the press reports
``writing \ldots'' immediately and \param{written = true} arrives on the
completion re-solve, so a slow network share never freezes the canvas.
The component persists \param{\_lastPath}/\param{\_lastInfo}
across sessions so re-opening the file still shows the last write's
outcome.

% =====================================================================
\section{Robot selection and simulation: Panel 4 Robot}
% =====================================================================

\subsection{Robot Definition}

\emph{``Select a robot or Cartesian machine profile for kinematics,
visualization and program generation.''} Not Heavy. Emits a stable,
versioned string, \texttt{iamlab-robot:v1:<profile-id>}, which every
robot-dependent component accepts on a single \param{robot} input.
\param{RobotCatalog.Resolve} also accepts canonical ids and short
aliases (\texttt{"scara"}, \texttt{"none"}), so a
plain text panel works as well as the component. An unknown definition
raises a clear error instead of silently resolving to the wrong
machine.

\begin{longtable}{@{}p{4.2cm}p{2.2cm}p{2.6cm}p{3.6cm}@{}}
\toprule
\textbf{Profile id} & \textbf{Mechanism} & \textbf{Program} & \textbf{Direct control} \\
\midrule
\endhead
\param{none} & none & RepRap G-code & \\
\param{generic-cartesian} & Cartesian & RepRap G-code & \\
\param{generic-duet-cartesian} & Cartesian & RepRap G-code & Duet (7 Print) \\
\param{potterbot-scara-v4} & SCARA & RepRap G-code & Duet (7 Print) \\
\bottomrule
\end{longtable}

Adding a robot is one \param{RobotCatalog} entry, plus a solver in
\param{Core/Robot/} only if it needs IK, plus a mesh provider only if it
should be drawn (\emph{ARCHI.md} \S 12.1). A Cartesian machine needs
neither.

\subsection{Robot Simulation}

\emph{``Visualize G-code/toolpaths and, when a Robot Definition supplies
kinematics, simulate and validate the selected robot.''} Replaces the
pre-2.0 \param{Robot} component (hidden but still solving in
saved documents). The numerical
solvers were not changed by the unification: the dispatch on the wired
profile selects which validator runs.

\begin{longtable}{@{}p{3.2cm}p{2.0cm}p{2.0cm}p{6.4cm}@{}}
\toprule
\textbf{Parameter} & \textbf{Unit} & \textbf{Default} & \textbf{Description} \\
\midrule
\endhead
\param{robot} & definition & & Empty/\param{none} = model-free toolpath and extrusion viewer. \\
\param{gcode} / \param{toolpath} & & required & Program (one string or lines), or curves used when G-code is empty. \\
\param{sim} & 0 to 1 & 1 & Scrub through the accumulated print. \\
\param{showRobot} / \param{showExtrusion} & bool & true & Viewport preview toggles. \\
\param{width} / \param{height} & mm & 10 / 0 & Layer preview cross-section (\param{0} = infer from layer metadata). \\
\param{bed} & definition & & Optional \param{Print Bed}; the explicit bed inputs override it when wired. \\
\bottomrule
\end{longtable}

\textbf{Groups:} \textbf{bed} (\param{bedWidth} 1000, \param{bedDepth}
700, \param{bedCenter}: visualization extents); \textbf{robotopts}
(\param{baseOffset}: bed/toolpath origin in robot-base coordinates;
\param{armMode}: SCARA arm mode used when the path is not anchored by
homing: a path with a \texttt{G28} takes its arm mode from the homing
endstops of the machine).

\textbf{Outputs:} \param{joints} (Number tree, one branch per joint:
empty without kinematics), \param{arm}/\param{tcp}/\param{bed} (preview
geometry), \param{reachOK}, \param{problems} (flagged target
positions), \param{messages}, \param{info}. Any target outside the limits of the selected machine also raises a hard red runtime
\textbf{Error} (``targets are outside robot limits'').
\textbf{Never bypass or silence that error}: if it fires, the fix is to
change the print (geometry, placement, or speeds), not to hide the
message. Deposition and travel draw in different colours (RoyalBlue
layer vs.\ thinner ochre travel), in both display modes.

\subsection{SCARA profile: ``Potterbot Scara v4''}

The SCARA validator walks the program with
\param{ScaraKinematics}/\param{RobotChecks} against the shipped
Potterbot profile: link lengths $L1 = 549.39$\,mm / $L2 = 458.62$\,mm
(firmware \texttt{M669 P/D}), reach annulus 235\,mm to $L1{+}L2$, $z$
range 0 to 1143\,mm, \param{maxFeed} 150\,mm/s, \param{maxJointSpeed}
120\iamdeg/s, joint ranges $J1 = -311.1$\iamdeg{} to $49.1$\iamdeg,
$J3 = \pm 151.93$\iamdeg{} (the \emph{relative} elbow angle: firmware
naming), $J2 = 0$ to 1030\,mm (the prismatic Z-carriage), and a
near-singularity jerk flag. It also models the real Cartesian workspace
and emulates the arm-mode of the firmware and joint-wrap choices, so a
toolpath that validates is one the machine will follow.

J1's firmware range spans 360.2\iamdeg{} but wraps at the $\pm 180$\iamdeg{}
boundary: angle comparisons for range-checking accept $\pm 360$\iamdeg{}
equivalents (fixed 2026-07-14 after a real $-200$\iamdeg{} pose, in range,
normalised to $+160$\iamdeg{} and incorrectly flagged as unsafe).

% =====================================================================
\section{Analyze components: Panel 5 Analyze}
% =====================================================================

\subsection{Overhang \heavy{}}

\emph{``Per-layer / per-segment overhang analysis for printability
assessment.''}

\begin{longtable}{@{}p{3.2cm}p{2.0cm}p{2.0cm}p{6.4cm}@{}}
\toprule
\textbf{Parameter} & \textbf{Unit} & \textbf{Default} & \textbf{Description} \\
\midrule
\endhead
\param{contours} & tree & required & From Slice or Place. \\
\param{warnDeg} & \iamdeg{} & 25 & Warning threshold: \textbf{placeholder, not calibrated}. \\
\param{failDeg} & \iamdeg{} & 45 & Failure threshold, clamped $\geq$ \param{warnDeg}; \textbf{placeholder}. \\
\param{tol} & mm & 0.5 & Segment tolerance. \\
\param{colorMode} & & 0 & 0 = three classes, 1 = continuous magma ramp from 0 to \param{failDeg}. \\
\param{global} & & off & Off = local reference (angle from the layer normal, the stagger); on = global reference (angle from world Z, the lean against gravity). Picks which reading the thresholds, colours and legend use. \\
\param{csvPath} & path & & Export group. \\
\bottomrule
\end{longtable}

\textbf{Outputs:} \param{angles} (the \textbf{primary} metric: each
segment's angle vs. the \emph{local} layer-surface normal, estimated via
a windowed Newell-normal fit over an approximately 15-point window around
the segment; for planar layers this is only the world-Z angle, but for
non-planar (conical) layers it measures actual layer-on-layer
stagger, which is what limits printability there),
\param{anglesWorld} (the gravity-relative angle vs world Z: switch
\param{global} on to classify by it if gravity load is considered the
dominant failure mode for a given mix rather than stagger),
\param{summary} (CSV text), \param{info}, \param{legend} (a wire for
Legend) and \param{colors}. Colour-coded viewport preview in
magma hues, three classes or a continuous ramp.

\textbf{\param{warnDeg}/\param{failDeg} are explicitly uncalibrated
placeholders}: published 3D-concrete-printing studies put safe
unsupported wall inclination around 15 to 35\iamdeg, and the shipped defaults
sit inside that envelope, but they are not measured values for any
specific material. Calibrate via a staircase-cone test print and set the
measured angles as plain inputs (they are not constants).

\subsection{Print Stats}

\emph{``Lengths, extruded volume/mass, time estimate and per-layer
breakdown of a print job.''} Not Heavy.

\begin{longtable}{@{}p{3.2cm}p{2.0cm}p{2.0cm}p{6.4cm}@{}}
\toprule
\textbf{Parameter} & \textbf{Unit} & \textbf{Default} & \textbf{Description} \\
\midrule
\endhead
\param{gcode} / \param{toolpath} & & required & Source. \\
\param{robot} & definition & & Optional; enables the joint-speed-aware time model (see below). \\
\param{width} & mm & 4 & Layer width, for the geometric volume fallback. \\
\param{height} & mm & 0 & \param{0} = auto-derive from Z-steps. \\
\param{density} & g/cm$^3$ & 2.0 & Material density (clay $\approx 2.0$). \\
\param{ramDiameter} & mm & 94 & \textbf{If set and the G-code has E values, volume is computed exactly from E}; otherwise falls back to the geometric width$\times$height estimate. \\
\param{printFeed} / \param{travelFeed} & mm/s & 25 / 60 & Feeds group. \\
\param{accel} & mm/s$^2$ & 0 (off) & Non-zero enables a trapezoidal accel-aware time model instead of the naive distance/feed estimate. \\
\bottomrule
\end{longtable}

\textbf{Outputs:} \param{report} (formatted text), \param{timeMin},
\param{volumeCm3}, \param{massG}, \param{extrudeLen}, \param{travelLen},
\param{layers}, \param{perLayer} (tree of \texttt{[z, length, time]} per
layer).

\textbf{Robot-aware time (2.1.0).} With a kinematic \param{robot}
definition wired, each segment's feed is additionally capped by the joint-speed limits of the machine, inverting the same check
\param{RobotChecks} enforces:
\[
feed \le \min(\text{requested},\ MaxFeed,\ MaxJointSpeed \cdot dist/|\Delta J_i|)
\]
Near the reach boundary or the dead zone a small Cartesian step demands
a large joint move, the cap collapses, and the segment correctly takes
longer. The joint-limited time is reported \emph{alongside} the
Cartesian estimate with a percentage delta, and \param{timeMin}
publishes the most constrained model available (robot-limited, else
accel-aware, else the naive lower bound). \param{robot = none} and
non-kinematic profiles leave every existing path bit-identical.

% =====================================================================
\section{Calibration workflow: Panel 6 Calibrate}
% =====================================================================

A structured way to find good print-parameter settings for a new
material or machine instead of guessing, built around four components
used in sequence:

\begin{verbatim}
DoE Design --> (print + measure) --> RSM Evaluate --> Desirability
        ^
        |
Test Geometry (sample shapes for the print runs)
\end{verbatim}

\subsection{DoE Design}

Generates a run plan for a chosen set of factors (parameter names,
low/high bounds) and a design type: \textbf{DSD} (definitive screening
design, 4 to 12 factors: good for a first pass to find which factors
matter), \textbf{BBD} (Box to Behnken, 3 to 7 factors, 15/27/46/54/62 runs: supports a full quadratic response surface), \textbf{Taguchi}
(L4/L8/L9/L12/L16/L18/L27 orthogonal arrays), \textbf{FullFact2} (full
2-level factorial, 1 to 10 factors), \textbf{FracFact2} (fractional
factorial, 4 to 9 factors, named generator sets with a reported
resolution).

\begin{longtable}{@{}p{3.2cm}p{9.3cm}@{}}
\toprule
\textbf{Parameter} & \textbf{Description} \\
\midrule
\endhead
\param{names}, \param{lows}, \param{highs} & Factor names and bounds. \\
\param{design} & Default \texttt{"BBD"}. \\
\param{responses} & Response names (creates empty columns to fill in). \\
\param{centres}, \param{centerPts}, \param{levels}, \param{seed}, \param{fracRuns} & Advanced group. \\
\param{sigma}, \param{effect}, \param{sims}, \ldots & Power group (2.3.0): Monte-Carlo power analysis: ``can this 15-run plan detect a 0.3\,mm effect'' answered before any clay is mixed, with a warning under the conventional 0.8. \\
\param{xlsxPath}, \param{csvPath} & Export group. \\
\bottomrule
\end{longtable}

\textbf{Outputs:} \param{runs} (actual-units tree), \param{coded} ($-1$
to $+1$ tree), \param{factorNames} (feeds \param{RSM Evaluate}),
\param{runLabels}, \param{table}, \param{power} (when the power group is
active), \param{info}.

\subsection{Printing and measuring the runs}

Use \textbf{Test Geometry} to generate a consistent sample coupon
per run, slice/place/G-code it with the settings from that DoE row,
print it, and measure the chosen responses by hand (e.g. dimensional
accuracy, surface finish score, collapse/overhang behaviour, cycle
time). Right-click dropdown shapes: Cylinder, Vase, TwistedVase,
Frustum, Wall, Box.

\begin{longtable}{@{}p{3.2cm}p{2.0cm}p{2.0cm}p{6.4cm}@{}}
\toprule
\textbf{Parameter} & \textbf{Unit} & \textbf{Default} & \textbf{Description} \\
\midrule
\endhead
\param{plane} & & & Placement group. \\
\param{height} & mm & 100 & Overall height. \\
\param{radius} & mm & 40 & Overall radius. \\
\param{waveAmp}, \param{waveCount}, \param{twistDeg}, \param{sides}, \param{topRatio} & mixed & & Profile group (shape-dependent). \\
\param{length}, \param{width} & mm & & Wall/box group. \\
\bottomrule
\end{longtable}

\textbf{Outputs:} \param{geometry} (feeds \param{Slice}),
\param{sections}, \param{info}.

\subsection{RSM Evaluate \heavy{}}

Fits a response-surface model (Linear / 2FI / Quadratic) per response
via ordinary least squares (QR decomposition), reporting coefficients
with standard error/t/p-values, R\textsuperscript{2}/adjusted-R\textsuperscript{2}/predicted-R\textsuperscript{2},
RMSE, and a full ANOVA including a lack-of-fit vs.\ pure-error split.

\begin{longtable}{@{}p{3.2cm}p{9.3cm}@{}}
\toprule
\textbf{Parameter} & \textbf{Description} \\
\midrule
\endhead
\param{tablePath} & Filled-in measurements file. \\
\param{factorNames} & From \param{DoE Design}. \\
\param{responseNames} & Auto-detected if left blank. \\
\param{model} & Linear / 2FI / Quadratic. \\
\param{lows}, \param{highs}, \param{centres} & Ranges group. \\
\param{X}, \param{Y} & Direct-data group (bypasses the file if data already lives in Grasshopper). \\
\param{xlsxReport} & Export group. \\
\bottomrule
\end{longtable}

\textbf{Outputs:} \param{models} (feeds \param{Desirability}),
\param{report}, \param{equations}, \param{predicted}, \param{json},
\param{diag} (2.3.0: leverage, studentized residuals, Cook's distance
per run, VIF per term, with warnings on Cook's $D > 1$ / VIF $> 10$),
\param{info}. \textbf{Watch for a weak-fit warning (adjusted-R\textsuperscript{2}
$< 0.5$)}: a weak model here means the downstream Desirability
recommendation is only as good as a poor fit, not a reliable optimum.

\textbf{File formats.} \param{tablePath} accepts Excel or CSV. CSVs are
detected per file as English (\texttt{1.5}) or German (\texttt{1,5})
number format; mixed conventions, thousands separators or malformed
quoting raise a \param{CsvFormatException} naming the line, independent
of the Windows locale: never a silently wrong number.

\subsection{Desirability \heavy{}}

Takes the fitted \param{models} plus a \textbf{goal string per
response} (a small DSL: \texttt{max}, \texttt{min}, \texttt{target=7 in
6..8}, \texttt{range=L..H}, with optional weight \texttt{w=} and shape
\texttt{s=} modifiers) and searches for the factor settings that best
satisfy all goals simultaneously, using Derringer to Suich weighted
geometric-mean desirability, a grid search (capped at 500{,}000 points)
plus Nelder to Mead refinement for precision.

\begin{longtable}{@{}p{3.2cm}p{9.3cm}@{}}
\toprule
\textbf{Parameter} & \textbf{Description} \\
\midrule
\endhead
\param{models} & From \param{RSM Evaluate}. \\
\param{goals} & Goal DSL strings, one per response. \\
\param{factorPriors} & Priors group: biases the search toward preferred regions. \\
\param{lows}, \param{highs}, \param{gridPerDim} & Bounds group; \param{gridPerDim} default 15. \\
\param{sweepResponse}, \param{sweepValues}, \param{sweepBand} & Sweep group: reproduces a ``recipe per target value'' style campaign. \\
\param{xlsxPath} & Export group. \\
\bottomrule
\end{longtable}

\textbf{Outputs:} \param{settings} (recommended actual-unit factor
values), \param{D} (overall desirability, 0 to 1), \param{predicted},
\param{report}, \param{alternatives} (\emph{local optima of the
scalarised D}, not a Pareto set), \param{sweep}, \param{interval}
(2.3.0: 95\,\% prediction/confidence intervals; alternatives whose
predictions sit inside the interval of the recommendation are marked
\texttt{$\approx$best} rather than ranked), and: with the optional
\textbf{pareto} group enabled: \param{pareto} (a non-dominated front
over the individual goal desirabilities, with a 2-objective scatter in
the report; 3+ objectives get a note instead of a misleading 2D
projection).

A Graph-Mapper-style \textbf{goal-curve strip} (prototype, 2.3.0)
appears below the desirability profile after the first solve: drag the
ramp ends, target, or shape handle instead of typing goal text. Edits
are stored as per-response overrides of the wired goal string (never
modifying it), survive save/reopen, and reset with the \texttt{X}
control; weight stays text-only and there is no undo.

% =====================================================================
\section{Utility components: Panel 1 Slice}
% =====================================================================

\subsection{Graph Mapper}

\emph{``Remap numbers through an editable curve: drag the grips on the
canvas.''} The stock Graph Mapper's job with improved mechanics: pick a
curve family (including custom point curves you own), drag grips
directly on the component, and feed the remapped stream into any
numeric input: a flow or speed gradient over layers is the typical
use. Styling follows the Grasshopper canvas skin.

% =====================================================================
\section{Rationale: pressure and pump}
% =====================================================================

\subsection{Pump extrusion}

The \param{pump} group on \param{G-code Generator} computes the
volumetric flow the pump needs to deliver:
\[
Q = h \cdot w \cdot v \cdot \eta
\]
(layer height $\times$ layer width $\times$ print speed $\times$
\param{slipFactor}), converted from mm$^3$/s to l/min, using
\param{firstLayerSpeedFactor} for the first layer. Example from the test
suite: $h=10$\,mm, $w=30$\,mm, $v=50$\,mm/s $\rightarrow$ \textbf{0.9
l/min}. \param{slipFactor} (default 1.0) is a calibration term
back-calculated after a real print by comparing planned vs.\ actually
deposited volume: typically $\geq 1.0$.

\textbf{This group does not touch the emitted G-code.} The E-values in
the file remain exactly the ram/plunger volumetric model's output,
unchanged, regardless of whether the pump group is on. This is
deliberate: a screw pump is driven by its own speed interface (typically
an analog 0 to 10\,V input), not the motion controller's E-axis, so there is nothing
meaningful to write into the G-code without knowing the specific control
wiring: leaving E intact is harmless if pump-driven and still correct
if ram-driven. (If a controller does drive the pump: e.g.\ a Duet
\texttt{M950}/\texttt{M42} PWM mapping: a flavor-specific control line
could be added; it is not currently implemented.)

\textbf{Calibration is measured, never hardcoded.} \param{pumpCalibCsv}
is a \texttt{setting, flow\_lpm} table; the lookup is
\textbf{piecewise-linear interpolation only}: it will not extrapolate
beyond the measured range (a warning and no recommendation, rather than
a guess, results outside that range). This is an explicit policy
decision: no pump's rated range is built in, because manufacturer
figures differ between models and listings, and even correct numbers
would not account for a specific material's rheology, pump wear, or
hose setup. The measured table's flow range is the only range. The shipped CSV is a placeholder
full of \texttt{NEEDS\_INPUT} markers.

\paragraph{Bucket-test calibration protocol (approximately 15 minutes).}
\begin{enumerate}
  \item Mix material and prime the line.
  \item Choose 4 to 6 settings spanning the working range (e.g.\ 2, 4, 6, 8, 10\,V).
  \item Run each into a tared bucket for 60 seconds and weigh.
  \item Measure fresh material density separately with a 1-litre cup.
  \item Compute \(flow\_lpm = \frac{mass_g}{density_{g/ml} \times 1000 \times time_{min}}\).
  \item Enter rows in ascending setting order in \param{pumpCalibCsv}.
\end{enumerate}
Repeat whenever material/mix or hose length changes: this is not a
one-time calibration.

\subsection{Legacy calibration values (30\,mm ram machine)}

The only surviving \emph{field-measured} calibration for a smaller,
30\,mm-bore ram cartridge machine (as opposed to the current 94\,mm-ram
defaults), preserved from the deleted Python v1 toolkit's
\param{RamGcodeBody.py}:

\begin{longtable}{@{}lrr@{}}
\toprule
\textbf{Parameter} & \textbf{30\,mm-ram (legacy, measured)} & \textbf{94\,mm-ram (current default)} \\
\midrule
\endhead
\param{layerHeight} & 3.02\,mm & 2.0\,mm \\
\param{width} & 7.0\,mm & 4.0\,mm \\
\param{pressureCoeff} & 1.3 & 1.0 \\
\param{printSpeed} & 43.0\,mm/s & 20.0\,mm/s \\
\param{travelSpeed} & 40.0\,mm/s & 60.0\,mm/s \\
\param{retractSpeed} & 10.0\,mm/s & 2.0\,mm/s \\
\param{pressureOffset} & 1.0\,mm & 0.1\,mm \\
\param{zHop} & 2.0\,mm & 0.0\,mm \\
\bottomrule
\end{longtable}

If printing on a 30\,mm-ram machine, start from the legacy column: it
was deliberately rescaled for the 94\,mm defaults, so do not ``correct''
the current defaults back toward these unless the hardware actually
matches the 30\,mm case. Note this table reflects the historical Python
prototype's defaults (\param{width} 4.0\,mm, \param{printSpeed}
20\,mm/s), which differ from the current defaults of the plugin listed
in \S 5 (\param{width} 10.0\,mm, \param{printSpeed} 43\,mm/s): the defaults of the plugin have since been retuned independently of this legacy
record.

\subsection{A real historical bug worth knowing about (ArcWelder)}

Until fixed on 2026-07-14, the ArcWelder CLI integration passed a path
tolerance as a bare number (\texttt{5.0}) where the executable actually
expects a fraction in $(0,1)$ (\texttt{0.05} = 5\%): the CLI rejected
every call and the code \textbf{silently fell back to unwelded (plain
G1) output} with no visible error. Any G-code generated by
\param{arcwelder} from before this fix was never actually welded despite
reporting success; the fix now also surfaces the CLI's exit code and
stderr on failure instead of silently succeeding.

% =====================================================================
\section{Known limits and explicit placeholders}
% =====================================================================

The following are shipped defaults that make the tool work out of the
box, \textbf{not} calibrated values for any particular material or
machine:

\begin{itemize}
  \item \param{Overhang}'s \param{warnDeg}/\param{failDeg} (25\iamdeg/45\iamdeg).
  \item The pump group's shipped \param{pumpCalibCsv}: placeholder, \texttt{NEEDS\_INPUT}.
  \item The \param{ecomp} group's \param{eCompRefMul} (1.0 = no-op): uncalibrated.
\end{itemize}

Other known limits:

\begin{itemize}
  \item An \textbf{unwired} \param{robot} input on the G-code generator keeps the legacy Potterbot coordinate frame (mirror about $y = -x$) for backwards compatibility: wire a definition to get the identity frame a normal Cartesian machine expects. The active frame is always visible in \param{info}.
  \item Extrusion volume off-planar (adaptive/conical modes) is only nominal unless \param{varh} is enabled on the G-code component.
  \item Time estimates are a lower bound unless \param{Print Stats}'s \param{accel} input is set non-zero (opt-in trapezoidal model) or a kinematic robot definition is wired (joint-limited model).
  \item Unwired Duet control components (7 Print) default to the lab's Potterbot SCARA; the SCARA calibration components are SCARA-only.
  \item A known pre-existing bug in wall offsetting around hole orientation is logged but not yet fixed.
\end{itemize}

% =====================================================================
\section{Troubleshooting}
% =====================================================================

\begin{longtable}{@{}p{4.6cm}p{7.8cm}@{}}
\toprule
\textbf{Symptom} & \textbf{Check} \\
\midrule
\endhead
Heavy component shows stale/empty result & Click its compute button: it does not auto-recompute. \\
\param{spacing} throws a hard error & Must be $> 0$. \\
Volume/time implausible & Does \param{layerHeight} match between Slice and G-code Generator? \\
Conical/adaptive volume looks off & Enable the \param{varh} group on G-code Generator. \\
Robot Simulation shows a red limits error & Do not bypass: fix geometry/placement/speed, re-simulate. \\
G-code looks mirrored on a Cartesian machine & Wire a Robot Definition: unwired keeps the legacy Potterbot frame (check \texttt{frame:} in \param{info}). \\
Part lands somewhere odd after changing document units & Heavy caches are invalidated on a unit change: press Compute again; mm-labelled inputs stay physical mm. \\
RSM/pump CSV refused & The named line mixes English/German number formats (or has broken quoting): export one consistent convention. \\
Pump group changes nothing in G-code & Correct: pump only calculates a setting, never touches E. \\
ArcWelder output looks unwelded & Confirm the CLI actually ran (info/exit code): was a real historical bug, fixed 2026-07-14. \\
Spiralize falls back / fails & A layer has $\neq 1$ closed contour: simplify the profile. \\
Save G-code never writes & \param{write} needs a Button (one-shot), not a persistent Toggle. \\
Save reports ``writing \ldots'' but no \param{written = true} & Writes are asynchronous: the confirmation arrives on the follow-up solve; a failure surfaces as a warning instead. \\
RSM Evaluate fit looks unreliable & Check adjusted-R\textsuperscript{2}: below 0.5 is a flagged weak fit. \\
\bottomrule
\end{longtable}

% =====================================================================
\section{Testing and provenance}
% =====================================================================

\param{Lugworm.Core.Tests} (980+ tests as of 2.6.0) runs the
algorithmic core independent of Grasshopper/Rhino: arc fitting,
offset/clipping, DoE design generation and inference, G-code
parsing/restart, geometry validation, robot
IK/limit checks (SCARA and Cartesian gates), machine-fingerprint
trust, CSV format detection, SemVer/update check, and the
pump/variable-flow/slope-compensation math each have dedicated test
files. If \param{Lugworm.Core} is modified, run this suite before
trusting new behaviour in Grasshopper.

\param{docs/DECISION\_LOG.md} in the project root records every
non-obvious design decision (the ambiguity, the choice, the reason, and
how to change it) across all development phases: consult it before
assuming a behaviour is a bug rather than a deliberate, documented
choice. \param{docs/AUDIT\_PHASE0*.md} and \param{docs/RUN\_REPORT*.md}
record point-in-time snapshots of what was fixed or added in each
development pass, including the ArcWelder and robot-fingerprint fixes
described in section 8.3 and section 2.1.

\end{document}
